\chapter{Analyse des Besoins}
\label{chap:analyse-besoins}

\section{Introduction}
Ce chapitre a pour objectif d'identifier et de structurer les besoins de la plateforme. Il presente le perimetre fonctionnel, les acteurs, les besoins fonctionnels et non fonctionnels, ainsi que les contraintes du projet.

\section{Identification des acteurs}
Les principaux acteurs du systeme sont :
\begin{itemize}[label=--]
    \item \textbf{Administrateur RH} : gere les CV, lance les traitements d'optimisation et de transformation, suit l'avancement global.
    \item \textbf{Gestionnaire Recrutement} : consulte les resultats, valide les sorties et declenche les notifications des candidats.
\end{itemize}

\section{Besoins fonctionnels}
Les besoins fonctionnels de la plateforme sont organises en modules.

\subsection{Module de gestion des CV}
\begin{itemize}[label=--]
    \item importer un ou plusieurs CV,
    \item visualiser la liste des CV importes,
    \item consulter les metadonnees d'un CV (identifiant, date, statut),
    \item supprimer un CV non pertinent.
\end{itemize}

\subsection{Module d'extraction et d'analyse}
\begin{itemize}[label=--]
    \item extraire automatiquement les informations principales (etat civil, formation, experiences, competences),
    \item structurer les informations extraites dans un format exploitable,
    \item signaler les donnees manquantes ou incoherentes.
\end{itemize}

\subsection{Module d'optimisation intelligente}
\begin{itemize}[label=--]
    \item ameliorer la qualite redactionnelle du contenu,
    \item reorganiser les sections du CV selon les bonnes pratiques,
    \item adapter le contenu au template cible tout en conservant le sens des informations.
\end{itemize}

\subsection{Module de transformation et generation}
\begin{itemize}[label=--]
    \item selectionner un template de branding,
    \item previsualiser le CV transforme avant export,
    \item generer le CV final au format DOCX.
\end{itemize}

\subsection{Module de notification}
\begin{itemize}[label=--]
    \item selectionner les candidats concernes,
    \item envoyer des notifications automatisees par email,
    \item tracer l'etat des envois (envoye, echec, en attente).
\end{itemize}

\section{Besoins non fonctionnels}
La qualite globale de la plateforme doit respecter les exigences suivantes :
\begin{itemize}[label=--]
    \item \textbf{Performance} : temps de reponse acceptable pour l'import, l'analyse et la generation.
    \item \textbf{Fiabilite} : execution stable des traitements avec gestion des erreurs.
    \item \textbf{Securite} : protection des donnees personnelles des candidats et controle d'acces.
    \item \textbf{Ergonomie} : interface simple, claire et intuitive pour les equipes RH.
    \item \textbf{Maintenabilite} : architecture modulaire facilitant les evolutions.
    \item \textbf{Scalabilite} : capacite a traiter un volume croissant de CV.
\end{itemize}

\section{Contraintes du projet}
Le projet est encadre par un ensemble de contraintes :
\begin{itemize}[label=--]
    \item contrainte temporelle liee a la periode du PFE,
    \item respect du contexte technique de l'entreprise d'accueil,
    \item prise en compte de la qualite des donnees d'entree (formats de CV variables),
    \item necessite de produire un resultat exploitable en environnement reel.
\end{itemize}

\section{Specifications des cas d'utilisation principaux}
\subsection{Cas d'utilisation UC1 : Importer un CV}
\textbf{Acteur principal :} Administrateur RH

\textbf{Precondition :} l'utilisateur est authentifie et dispose des droits necessaires.

\textbf{Scenario nominal :}
\begin{enumerate}
    \item l'utilisateur accede au module d'import,
    \item il selectionne un fichier CV,
    \item le systeme verifie le format,
    \item le systeme enregistre le fichier et confirme l'import.
\end{enumerate}

\textbf{Postcondition :} le CV est disponible dans la liste des profils a traiter.

\subsection{Cas d'utilisation UC2 : Transformer et generer un CV}
\textbf{Acteur principal :} Gestionnaire Recrutement

\textbf{Precondition :} un CV valide est disponible dans le systeme.

\textbf{Scenario nominal :}
\begin{enumerate}
    \item l'utilisateur choisit un profil,
    \item il selectionne un template,
    \item le systeme optimise et restructure le contenu,
    \item le systeme affiche un apercu,
    \item l'utilisateur valide et lance la generation DOCX.
\end{enumerate}

\textbf{Postcondition :} le document final est genere et telechargeable.

\subsection{Cas d'utilisation UC3 : Notifier un candidat}
\textbf{Acteur principal :} Gestionnaire Recrutement

\textbf{Precondition :} le candidat est identifie et son email est disponible.

\textbf{Scenario nominal :}
\begin{enumerate}
    \item l'utilisateur selectionne un ou plusieurs candidats,
    \item il choisit un modele de message,
    \item le systeme envoie les notifications,
    \item le systeme enregistre le statut d'envoi.
\end{enumerate}

\textbf{Postcondition :} les candidats recoivent les notifications et le suivi est journalise.

\section{Priorisation des besoins}
\begin{table}[H]
\centering
\begin{tabular}{p{6cm}p{3cm}p{3cm}}
\toprule
\textbf{Besoin} & \textbf{Priorite} & \textbf{Iteration cible} \\
\midrule
Import et gestion des CV & Haute & Iteration 1 \\
Extraction des informations & Haute & Iteration 1 \\
Optimisation intelligente & Haute & Iteration 2 \\
Transformation via templates & Haute & Iteration 2 \\
Generation DOCX & Haute & Iteration 2 \\
Notification email & Moyenne & Iteration 3 \\
Tableau de suivi avance & Moyenne & Iteration 3 \\
\bottomrule
\end{tabular}
\caption{Priorisation des besoins fonctionnels}
\label{tab:priorisation-besoins}
\end{table}

\section{Conclusion}
Ce chapitre a formalise les besoins de la plateforme, en distinguant les exigences fonctionnelles et non fonctionnelles. Ces specifications constituent la base de la conception detaillee presentee dans le chapitre suivant.
